% !PRINTED 388
\weberpage{388}

%<<CONT
and it therefore suffices to denote the whole continued fraction by the period, say thus:

\weq{10}{\omega = [\alpha_0, \alpha_1 \ldots \alpha_{\nu-1}].}

This theorem is evidently equivalent to saying that the closing number $\omega_\nu$ of the continued fraction (9) is identical with $\omega$ itself, and in this also lies the proof of the assertion. For the closing numbers $\omega_\nu$, being equivalent numbers, all belong to the same discriminant, and hence by 4. the number of all possible closing numbers is only a finite one. In the sequence of closing numbers $\omega, \omega_1, \omega_2 \ldots$ one that has already occurred must therefore at some point appear for the second time. If $\omega_\nu$ is the first that appears for the second time, then $\omega_\nu = \omega$ must hold, since otherwise by 3. \S.\,124 the number $\omega_{\nu-1}$ would also appear for the second time.

Thus all the reduced numbers belonging to one discriminant can be arranged in periods:

\weq{11}{\omega, \omega_1, \omega_2 \ldots \omega_{\nu-1},}

so that, if
\[
\alpha_0, \alpha_1, \alpha_2 \ldots \alpha_{\nu-1}
\]
are the greatest integers contained in them, we have
\[
\begin{array}{l}
\omega = [\alpha_0, \alpha_1 \ldots \alpha_{\nu-1}]\\
\omega_1 = [\alpha_1, \alpha_2 \ldots \alpha_0]\\
\cdots\cdots\cdots\cdots\cdots\cdots\cdots\cdots
\end{array}
\]
by which the continued fraction expansion is given for each of these numbers.

If the whole system of reduced numbers is not yet exhausted by \gesperrt{one} period, one forms a second period, and so on. Since by a number $\omega$ both the number preceding it in the period and the one following it are completely determined, two different periods never contain a common number.

Of these periods there now holds the theorem:

\gesperrt{6. Reduced numbers from the same period are equivalent; from different periods they are not equivalent.}

The first part of the assertion is clear from the outset, since reduced numbers of the same period are closing numbers of one another; the second part is proved by the theorem \S.\,121, VI, that equivalent numbers, when the continued fraction expansion is carried sufficiently
%CONT>>
